\chapter{Liderazgo, Carisma y Construcción de Equipos}\label{ch:leadership-charisma}

% Alcance: liderazgo basado en evidencia para fundadores — liderazgo transformacional (Judge y Piccolo),
%   tácticas de liderazgo carismático (CLTs) medibles y entrenables (Antonakis), selección de líderes
%   (aptitud mental general + responsabilidad; advertencia sobre IE), seguridad psicológica (Edmondson; Frazier et al.),
%   gobernanza del equipo fundador (descriptivo de Wasserman; marco NVT de Ensley).
%   NUEVO (Parte VII). Extiende \cref{ch:organizational-behavior} (incentivos/estructura/gobernanza)
%   desde el lado del liderazgo personal --- referencia cruzada, sin duplicar.
% Gancho de valor: la calidad del liderazgo impulsa la retención, la contratación y la consistencia en la ejecución -> ROIC.

El trabajo de un fundador no es ser el mejor colaborador individual en la sala. Es elevar el
desempeño de todos los demás. \Cref{ch:organizational-behavior} aborda las herramientas
estructurales disponibles para esa tarea: diseño de incentivos, gobernanza y arquitectura
organizacional. Este capítulo trabaja el mismo problema desde el lado personal --- qué hace, dice
y señala un líder en el flujo diario de decisiones. La conexión con la creación de valor es
directa: un liderazgo de alta calidad reduce la rotación involuntaria, comprime el tiempo
necesario para cubrir puestos directivos y sostiene la consistencia en la ejecución que convierte
un modelo operativo sólido en los números de ROIC que \cref{ch:corporate-finance-core} identifica
como el motor del valor empresarial. Nada de lo que sigue requiere dotes innatas. La base
empírica es inusualmente sólida: los modelos centrales están operacionalizados, probados y, en
dos casos, confirmados experimentalmente.

\section{Liderazgo Transformacional}\label{sec:transformational-leadership}
% complejidad: 2

¿Qué separa a los líderes que obtienen un desempeño competente de aquellos que elicitan un
esfuerzo extraordinario? El modelo de liderazgo transformacional --- desarrollado por
\textcite{bass1985leadership} y validado cuantitativamente por \textcite{judge2004transformational}
--- responde esa pregunta distinguiendo la mecánica de la \emph{transformación} de la mecánica de
la \emph{transacción}. El liderazgo transaccional está basado en el intercambio: el líder ofrece
recompensas por un buen desempeño y corrige las desviaciones del estándar. Es necesario para las
operaciones base, pero insuficiente para inspirar un esfuerzo que supere lo que especifica el
contrato. En una empresa en etapa temprana donde las recompensas financieras son escasas y el
futuro es incierto, la gestión transaccional está estructuralmente subpotenciada. El liderazgo
transformacional sustituye la estructura formal faltante por algo más difícil de replicar: un
sentido de propósito convincente que hace que los seguidores estén dispuestos a actuar más allá
de su interés inmediato.

El modelo operacionaliza el liderazgo transformacional a través de cuatro dimensiones conductuales,
ampliamente conocidas como las Cuatro~Is:
\begin{description}
  \item[Influencia idealizada] El líder actúa como un sólido modelo a seguir, manteniendo altos
    estándares de conducta ética. Los seguidores se identifican con el líder y desean emularlo.
    Aquí reside el componente de «carisma» comúnmente observado --- se divide en influencia
    idealizada atribuida (cómo perciben los seguidores al líder) e influencia idealizada
    conductual (lo que el líder hace en realidad).
  \item[Motivación inspiracional] El líder comunica altas expectativas a través de símbolos y
    apelaciones emocionales, enfocando el esfuerzo colectivo en metas que superan lo que el
    interés propio solo motivaría.
  \item[Estimulación intelectual] El líder desafía a los seguidores a cuestionar supuestos ---
    incluidos los propios y los del líder --- y a abordar los problemas desde nuevos ángulos,
    creando condiciones para el pensamiento independiente en lugar de la mera obediencia.
  \item[Consideración individualizada] El líder actúa como \emph{coach} y asesor, atendiendo a
    las necesidades particulares de desarrollo de cada seguidor en lugar de tratar al equipo como
    una audiencia uniforme.
\end{description}

La evidencia cuantitativa es sólida y específica. En un metaanálisis de \num{87} estudios
independientes que representan 626 correlaciones, \textcite{judge2004transformational} reportan
una correlación corregida de \(\rho = .58\) entre liderazgo transformacional y satisfacción
% fuente: judge2004transformational — rho corregido (correlación de puntaje verdadero) para
% LT→satisfacción laboral del seguidor = .58; LT→desempeño laboral del seguidor = .27 (Tabla 3, k=87)
laboral del seguidor, y de \(\rho = .27\) entre liderazgo transformacional y desempeño laboral
del seguidor. Estos no son hallazgos incidentales: el efecto sobre la satisfacción se encuentra
entre los más grandes documentados de manera confiable en la literatura del comportamiento
organizacional, y el efecto sobre el desempeño opera por encima y más allá de lo que los
incentivos transaccionales solos pueden producir. La implicación práctica es que un fundador que
invierte en conductas transformacionales no está ejerciendo una gestión blanda --- está desplegando
la palanca de influencia de mayor rendimiento que respalda la evidencia.

Bass y Riggio advierten explícitamente contra el liderazgo \emph{pseudo-transformacional}: líderes
que despliegan motivación inspiracional con fines narcisistas o de explotación fabrican la
apariencia superficial del modelo mientras invierten su intención. La autenticidad es estructural.
En el estudio longitudinal de \textcite{baum1998longitudinal}, la visión del fundador y su
comunicación no aceleraron directamente el crecimiento de la empresa; el efecto estuvo mediado
por el rigor y la consistencia de la transmisión verbal y escrita. Una visión que existe solo en
un \emph{pitch} no tiene palanca causal. La implicación para la ejecución es que el liderazgo
transformacional no es un rasgo de personalidad que se representa en reuniones generales --- es
un patrón sostenido de conducta a través de cada interacción, grande y pequeña.

\begin{example}[Comunicación de la visión en la reunión general]
La empresa SaaS tiene un \qty{42}{\percent} de margen de contribución y un ARR actual de
aproximadamente \money{4000000}. La CEO fundadora abre la reunión general trimestral no con la
cifra de ingresos, sino con una sola oración sobre el problema que la empresa existe para
resolver. Luego conecta el trabajo de cada equipo con una de las cuatro dimensiones
transformacionales: el esfuerzo de confiabilidad del equipo de ingeniería se vincula a la
influencia idealizada (la empresa hace lo que dice que hará); el proceso de descubrimiento del
equipo de ventas se vincula a la consideración individualizada de los clientes. Cierra
cuantificando la brecha entre el desempeño actual y la posibilidad que implica la misión --- no
para inducir ansiedad, sino para hacer que la distancia parezca salvable. Esta secuencia no es
accidental. \textcite{baum1998longitudinal} confirman que los atributos de la visión afectan el
crecimiento solo cuando se combinan con este tipo de comunicación consistente y estructurada; el
discurso es el mecanismo, no el adorno.
\end{example}

\begin{admonition}{Nota}
El efecto de desempeño de \(\rho = .27\) de \textcite{judge2004transformational} es un promedio
corregido entre estudios. El mayor valor práctico del modelo es diagnóstico: proporciona un
vocabulario para identificar qué dimensión está ausente cuando un equipo tiene un desempeño
inferior. Un equipo que ejecuta de manera confiable pero nunca mejora probablemente carece de
estimulación intelectual. Un equipo que genera ideas pero no puede ejecutar probablemente carece
de consideración individualizada y expectativas claras.
\end{admonition}

\section{El Carisma Es Entrenable}\label{sec:charisma-tactics}
% complejidad: 2

El carisma ha sido tratado durante la mayor parte de su historia intelectual como un don
cuasi-místico accesible a unos pocos elegidos. \textcite{antonakis2011charisma} desmanteló ese
supuesto experimentalmente. En un estudio controlado aleatorizado, los gerentes entrenados en
tácticas verbales y no verbales específicas fueron calificados significativamente más carismáticos
por compañeros de trabajo independientes tres meses después de la intervención, con un tamaño del
efecto estandarizado de aproximadamente
% fuente: antonakis2011charisma — d de Cohen ≈ 0.62 (RCT de entrenamiento en carisma, experimento de campo
% N=34 gerentes + experimento de laboratorio N=41 participantes MBA; calificaciones de pares a 3 meses)
\(d \approx 0.62\) --- un desplazamiento moderadamente grande según los estándares de la
psicología organizacional. La conclusión práctica es decisiva: la señalización carismática es una
habilidad aprendible, no un rasgo heredado al nacer, y los fundadores que la tratan como
imposible de enseñar están dejando inutilizada una herramienta de influencia sustancial.

Antonakis identifica nueve tácticas de liderazgo carismático (CLTs) verbales y tres no verbales.
Las tácticas verbales son:
\begin{itemize}
  \item \textbf{Metáfora y analogía} --- comprimir una visión compleja en una sola imagen concreta
    (\emph{«Estamos construyendo el sistema nervioso de internet»}).
  \item \textbf{Historias y anécdotas} --- anclar afirmaciones abstractas en una narrativa humana
    particular, que los oyentes procesan con menor resistencia cognitiva que el argumento
    declarativo.
  \item \textbf{Listas de tres elementos} --- una estructura retórica que crea la impresión de
    exhaustividad sin agotar a la audiencia (\emph{«Rápido, confiable y tuyo»}).
  \item \textbf{Preguntas retóricas} --- invitar a la audiencia a proporcionar la respuesta
    obvia, aumentando la adhesión mediante el acto de articulación.
  \item \textbf{Convicción moral} --- señalar que una decisión está impulsada por valores más que
    puramente por la ganancia, lo que eleva la confianza percibida en el líder.
  \item \textbf{Expresar el sentimiento colectivo} --- reconocer el estado emocional del equipo
    antes de pedirles que lo trasciendan.
  \item \textbf{Establecer altas expectativas} --- exigir excelencia en términos explícitos, lo
    que calibra lo que los seguidores creen que es alcanzable.
  \item \textbf{Expresar confianza} --- asegurar al equipo que el objetivo es alcanzable, lo que
    reduce el costo psicológico del esfuerzo bajo incertidumbre.
  \item \textbf{Contrastes y recursos retóricos} --- estructurar argumentos mediante la oposición
    (\emph{«No seguimos el mercado; lo creamos»}) para hacer que las posiciones sean memorables y
    claras.
\end{itemize}
Las tácticas no verbales --- gestos corporales deliberados, expresión facial alineada a la
valencia del mensaje y modulación vocal animada --- tienen un efecto independiente sobre el
carisma percibido y deben practicarse junto con el conjunto verbal.

El protocolo de entrenamiento de \textcite{antonakis2011charisma} sigue una secuencia de cinco
etapas: evaluación 360 grados de referencia, instrucción teórica en la taxonomía de CLTs,
redacción y revisión de discursos para incorporar las tácticas verbales, práctica grabada en
video para calibrar las tácticas no verbales, y reevaluación longitudinal a los tres meses. Los
fundadores que invierten en esta secuencia regresan siendo mensurablemente más persuasivos. La
dimensión de comunicación del entrenamiento conecta directamente con los marcos de
\cref{ch:storytelling-communication}, que construye la arquitectura narrativa que las CLTs están
diseñadas para entregar, y con \cref{ch:persuasion-influence}, que cubre los mecanismos
cognitivos y sociales subyacentes a la influencia de manera más amplia.

\begin{example}[Aplicación de CLTs en una reunión general del fundador]
La CEO de la SaaS está a punto de anunciar un giro significativo de producto. Estructura el
discurso utilizando cinco CLTs explícitamente. Abre con una \textbf{metáfora}: \emph{«Hemos
estado construyendo un caballo más rápido. Este giro es nuestra decisión de construir el
automóvil.»} Luego cuenta una \textbf{historia} sobre una llamada específica con un cliente en
la que los límites actuales del producto le costaron tres horas por semana. Formula una
\textbf{pregunta retórica}: \emph{«¿Es eso la empresa que estamos intentando ser?»} Declara su
\textbf{convicción moral} --- el giro es lo correcto para los clientes, no solo para el
crecimiento del ARR. Cierra con una \textbf{lista de tres elementos} de lo que el equipo
entregará en el próximo trimestre: el nuevo modelo de datos, el \emph{onboarding} rediseñado y
la primera integración empresarial. La misma información entregada como una actualización de
estado con viñetas habría sido procesada y luego olvidada. Entregada a través de CLTs, se
procesa e internaliza.
\end{example}

\begin{admonition}{Advertencia}
Antonakis advierte explícitamente que las CLTs pueden fallar cuando se usan en exceso o se
ejecutan de manera inauténtica. Un fundador que inserta metáforas en cada oración, o que expresa
convicción moral sobre temas donde los seguidores saben que la decisión fue comercial, socava el
valor de señal de las tácticas. Las CLTs funcionan porque son señales cognitivamente creíbles de
competencia y valores; su frecuencia debe permanecer coherente con esa afirmación de
credibilidad.
\end{admonition}

\section{Selección y Desarrollo de Líderes}\label{sec:leader-selection}
% complejidad: 2

¿Qué características individuales predicen realmente la eficacia del liderazgo? Esta es la
pregunta operativa que enfrenta un fundador cada vez que realiza una contratación directiva o
evalúa a un cofundador. La respuesta de la literatura de metaanálisis es lo suficientemente
precisa como para guiar un proceso de contratación: la aptitud mental general (GMA) y la
responsabilidad son los dos predictores más robustos y generalizables, las conductas
transformacionales añaden un valor incremental sustancial por encima de ellas, y la inteligencia
emocional (IE) de habilidad contribuye modestamente más allá de esa referencia.

La GMA --- la capacidad general de aprender, razonar y adaptarse --- es el predictor cognitivo
más fuerte del liderazgo. En un metaanálisis de \num{151} muestras independientes,
% fuente: judge2004intelligence — rho corregido GMA↔liderazgo = .27 (k=151; JAP 2004)
\textcite{judge2004intelligence} reportan una correlación corregida de \(\rho = .27\) entre GMA
y eficacia del liderazgo. El efecto no es grande en términos absolutos, pero es muy confiable en
distintos contextos y se mantiene después de la corrección por restricción de rango --- el
artefacto estadístico que tiende a suprimir las correlaciones de inteligencia en la investigación
de gestión, porque las organizaciones preseleccionan a los candidatos en umbrales cognitivos
mínimos.

En cuanto a la personalidad, la responsabilidad es el predictor más generalizable del desempeño
laboral en todos los grupos ocupacionales. En su metaanálisis fundamental de \num{117} estudios,
% fuente: barrick1991big — rho corregido Responsabilidad↔desempeño laboral ≈ .22 (k=117;
% JAP 1991); el predictor universal entre grupos ocupacionales
\textcite{barrick1991big} establecen una correlación corregida de aproximadamente \(\rho = .22\)
entre responsabilidad y desempeño laboral. La estabilidad emocional (bajo neuroticismo) sigue
como el segundo predictor más consistente de los Cinco Grandes (Big Five). La implicación
práctica para la contratación temprana es que un candidato que puntúa alto en responsabilidad ---
orientación a metas, autodisciplina, confiabilidad --- lleva una prima de desempeño documentada
que se mantiene independientemente del tipo de puesto o la industria, lo que lo convierte en la
señal de personalidad más defendible disponible para un fundador que construye un equipo pequeño
donde cada contratación importa desproporcionadamente.

Las conductas transformacionales, cuantificadas por \textcite{judge2004transformational} con
\(\rho = .58\) para la satisfacción del seguidor y \(\rho = .27\) para el desempeño del
seguidor, añaden una capa conductual significativa sobre estos predictores de diferencias
individuales. Una contratación que es cognitivamente capaz y responsable, pero que gestiona solo
a través de la transacción, deja un desempeño sustancial del equipo sobre la mesa. Seleccionar
para el potencial transformacional --- evaluándolo a través de entrevistas conductuales
estructuradas en lugar de autoreporte --- cierra esa brecha.

La inteligencia emocional (IE), en cambio, merece una advertencia calibrada.
% fuente: joseph2010emotional — IE de habilidad↔desempeño laboral r≈.19; validez incremental
% sobre GMA + Responsabilidad = ΔR²≈.002 (0.2%); metaanálisis de clase JAP
\textcite{joseph2010emotional} encuentran una correlación modesta de \(r \approx .19\) entre la
IE de habilidad y el desempeño laboral --- una relación que en gran medida desaparece una vez que
se controlan estadísticamente GMA y Responsabilidad, dejando solo una validez incremental
negligible. Los fundadores deben tratar los instrumentos de IE comercializados como señales
suplementarias en el mejor de los casos, y anclar su arquitectura de contratación en evaluaciones
validadas de GMA e inventarios estructurados de los Cinco Grandes en su lugar.

\begin{example}[Contratación del director de customer success]
La empresa SaaS está en aproximadamente \money{4000000} de ARR, \qty{42}{\percent} de margen
de contribución y un \emph{churn} mensual del \qty{0.65}{\percent}. La CEO está contratando a un
director de \emph{customer success} --- un puesto donde el desempeño afecta directamente la
retención y, por lo tanto, el CLV. Estructura el proceso en tres etapas. Primero, una evaluación
validada de aptitud cognitiva (filtro de GMA) para confirmar que el candidato puede manejar la
complejidad analítica y estratégica del puesto. Segundo, un inventario estructurado de los Cinco
Grandes para evaluar la responsabilidad y la estabilidad emocional --- los dos rasgos que
predicen de manera más robusta el desempeño sostenido. Tercero, una entrevista conductual que
cubre tres escenarios transformacionales: una ocasión en que el candidato entrenó a un miembro
del equipo con dificultades (consideración individualizada), una ocasión en que reencuadró un
contratiempo como una oportunidad de aprendizaje (estimulación intelectual), y una ocasión en
que alineó a un equipo en torno a un objetivo más allá de su interés inmediato (motivación
inspiracional). Los contratos con clientes por encima de \money{50000} requerirán la aprobación
de ambos cofundadores; el director de \emph{customer success} será propietario de todo lo que
esté por debajo de ese umbral. El proceso estructurado toma dos semanas y es sustancialmente más
predictivo que una serie de entrevistas no estructuradas --- y mucho más defendible cuando la
contratación no funciona.
\end{example}

\begin{admonition}{Nota}
Las evaluaciones de IE de autoreporte --- el tipo más ampliamente vendido a las organizaciones ---
miden la personalidad y la autopercepción, no la capacidad de procesamiento emocional. Son
susceptibles a la gestión deliberada de impresiones, lo que las hace poco confiables para
contrataciones de alto riesgo. El modelo de habilidad de la IE, medido a través de pruebas
basadas en el desempeño en lugar del autoreporte, se sustenta sobre una base más sólida, pero
incluso su correlación mejor validada con el desempeño laboral (\(r \approx .19\)) está
sustancialmente por debajo de lo que ofrecen GMA y Responsabilidad. Utilice los constructos que
respalda la evidencia.
\end{admonition}

\section{Seguridad Psicológica y Construcción del Equipo Fundador}\label{sec:psych-safety-team}
% complejidad: 3

La visión y el carisma de un líder proporcionan dirección y energía. Por sí solos, no protegen a
un equipo de la causa más común del fracaso en etapas tempranas: el conflicto interno y el
silencio que crea. \textcite{edmondson1999psychological} introdujo el concepto de seguridad
psicológica --- una creencia compartida del equipo de que asumir riesgos interpersonales es seguro
--- a través de un estudio de 51 equipos de trabajo que sigue siendo uno de los artículos más
citados en el comportamiento organizacional. El hallazgo fue preciso: la seguridad psicológica
media la relación entre la estructura del equipo y las conductas de aprendizaje del equipo. Los
equipos donde los miembros se sentían seguros para admitir errores, pedir ayuda y cuestionar
supuestos aprendían más rápido y tenían un mejor desempeño.

El efecto escala. En una revisión metaanalítica de \num{136} muestras independientes que
representan más de
% fuente: frazier2017psychsafety — metaanálisis k=136 muestras, N>22,000 individuos;
% la seguridad psicológica es un predictor dominante del aprendizaje del equipo, la voz y los resultados de desempeño
22,000 individuos, \textcite{frazier2017psychsafety} confirman la seguridad psicológica como un
predictor dominante de las conductas de aprendizaje del equipo, la voz del empleado y los
resultados de desempeño del equipo. La amplitud de la base de evidencia es importante: el efecto
se mantiene entre industrias, tamaños de equipo y contextos nacionales, posicionando la seguridad
psicológica no como un lujo cultural para equipos bien financiados, sino como un prerrequisito
operativo para cualquier equipo que necesite detectar y corregir errores rápidamente. Para una
empresa B2B SaaS cuyo riesgo central es el fallo invisible del producto y el colapso silencioso
del \emph{pipeline}, ese prerrequisito no es abstracto --- es el mecanismo que mantiene al CEO
informado antes de que los problemas se agraven.

El mecanismo no es comodidad ni amabilidad. La seguridad psicológica no es la ausencia de
responsabilidad; es la presencia de suficiente confianza interpersonal para que las conversaciones
de rendición de cuentas puedan ocurrir honestamente. \textcite{edmondson1999psychological} muestra
que la inclusividad del líder --- invitar activamente al disenso, reconocer la propia falibilidad
del líder y enmarcar el trabajo como un problema de aprendizaje en lugar de una prueba de
ejecución --- es el principal impulsor de la seguridad a nivel de equipo. Para un CEO fundador
que ya ha desarrollado las habilidades de señalización carismática descritas en
\cref{sec:charisma-tactics}, la inversión adicional relevante es crear momentos estructurales para
la retroalimentación honesta: postmórtem explícitos, reconocimiento público de los propios errores
del CEO y recompensa visible para el miembro del equipo que saca a la luz un problema
tempranamente en lugar de esperar a que se resuelva solo.

El problema de la estructura del equipo comienza antes de que cualquiera de estas inversiones
culturales sea posible. \textcite{wasserman2012founders} analizó datos de casi 10,000 fundadores y
encontró que los problemas con las personas --- no los fracasos de mercado o producto --- son la
causa principal del fracaso de las \emph{startup}. El sesenta y cinco por ciento de las empresas
de alto potencial que fracasan lo hacen debido a conflictos interpersonales entre cofundadores,
típicamente arraigados en visiones divergentes o expectativas desalineadas sobre roles y control.
Wasserman identifica tres dilemas que deben resolverse tempranamente, antes de que se consoliden
en conflicto.

Más allá de la evidencia descriptiva, el marco revisado por pares para comprender las dinámicas
de liderazgo en los equipos fundadores proviene de \textcite{ensley2006shared}, quien estudió los
equipos de alta dirección de nuevas empresas y encontró que el liderazgo compartido ---
distribuir la autoridad de decisión y
% fuente: ensley2006shared — el liderazgo compartido/vertical en los NVTs predice el desempeño
% de la nueva empresa; el liderazgo compartido añade 10%–15% de varianza incremental (Leadership Quarterly 2006)
la responsabilidad de liderazgo entre el equipo fundador --- predice el desempeño de la nueva
empresa con un \qty{10}{\percent} a \qty{15}{\percent} de varianza incremental por encima y más
allá de lo que contribuye solo un único líder vertical. El mecanismo es la autonomía específica
por dominio: el liderazgo compartido no significa consenso en cada decisión, sino más bien
distribución explícita de la autoridad última por área funcional (el cofundador técnico es
propietario de las decisiones de arquitectura; el cofundador comercial es propietario de la
estrategia de \emph{go-to-market}) con integración estructurada para la alineación de toda la
empresa. Este es el fundamento validado empíricamente para el protocolo de gobernanza que el
equipo fundador debe negociar antes de cualquier conversación de financiamiento externo.

\begin{description}
  \item[El Dilema de la Relación] Fundar con conexiones sociales previas proporciona comodidad
    inicial pero crea una renuencia estructural a abordar la incompetencia o las visiones
    divergentes más adelante. La familiaridad que facilita la comunicación temprana también
    suprime las conversaciones francas que se vuelven necesarias a medida que la empresa crece.
  \item[El Dilema de la Recompensa] Los repartos de capital iguales son la decisión temprana más
    común y la más frecuentemente lamentada. Un reparto igual no tiene en cuenta los diferentes
    niveles de compromiso continuo, contribución financiera y creación de valor futuro.
    \textcite{wasserman2012founders} muestra que una estructura de capital diferenciada negociada
    tempranamente --- y revisada explícitamente en los hitos de financiamiento --- está fuertemente
    asociada con la durabilidad de la relación entre cofundadores.
  \item[El Dilema del Control (Rico vs.\ Rey)] Los fundadores deben decidir eventualmente si
    maximizan el resultado financiero aceptando la dilución de inversores de primer nivel y
    ejecutivos experimentados (eligiendo ser «Ricos»), o si mantienen el control operativo y de
    gobernanza a costa de limitar el potencial de escala de la empresa (eligiendo ser «Reyes»).
    Ninguna opción es incorrecta, pero confundirlas sí lo es: los fundadores que quieren el
    resultado financiero del camino «Rico» mientras toman las decisiones de control del camino
    «Rey» crean crisis en cada evento de financiamiento. Este dilema conecta directamente con la
    mecánica de dilución de \cref{ch:fundraising-and-dilution}, donde la tabla de propiedad
    post-money hace visible la disyuntiva en números.
\end{description}

\begin{example}[Configuración del equipo fundador en la empresa SaaS]
Dos cofundadores lanzan la empresa SaaS. Ambos contribuyeron al producto inicial, pero uno
(el fundador técnico) continuará a tiempo completo desde el primer día; el otro (el fundador
comercial) está haciendo la transición desde un compromiso de consultoría durante seis meses.
Toman tres decisiones antes de escribir una línea de código de producción. Primero, dividen el
\emph{equity} en \qty{55}{\percent}/\qty{45}{\percent} en lugar de
\qty{50}{\percent}/\qty{50}{\percent}, ponderado al compromiso anterior a tiempo completo del
fundador técnico, con un calendario de consolidación de cuatro años y un período de un año de
consolidación mínima. Segundo, definen un protocolo de toma de decisiones coherente con el marco
de liderazgo compartido de \textcite{ensley2006shared}: las decisiones de arquitectura de
producto pertenecen al fundador técnico; los contratos con clientes por encima de \money{50000}
requieren ambos. Tercero, programan una revisión mensual entre cofundadores --- un punto fijo en
la agenda cuyo único propósito es sacar a la superficie cualquier cosa que alguno de los
fundadores haya estado renuente a plantear en el flujo ordinario del trabajo. Ese tercer
mecanismo es la infraestructura de seguridad psicológica para un equipo de dos personas. Su
existencia no previene el conflicto; asegura que el conflicto surja cuando aún es resoluble, en
lugar de cuando se ha agravado hasta convertirse en algo que la empresa no puede sobrevivir.
\end{example}

\begin{admonition}{Advertencia}
El hallazgo más consistente de \textcite{wasserman2012founders} es que los repartos de capital
decididos rápidamente y de manera igualitaria son un predictor principal de la disolución entre
cofundadores. La velocidad del acuerdo es la señal de alerta: un reparto concluido en una tarde
casi con certeza no ha trabajado las asimetrías en compromiso, tolerancia al riesgo y ambición
que se volverán decisivas doce a treinta y seis meses después. La intervención correcta no es
retrasar indefinidamente, sino usar una conversación estructurada --- que cubra el compromiso a
tiempo completo, las contribuciones financieras, la preferencia Rico-vs.-Rey y las salidas
aceptables --- antes de formalizar cualquier reparto. Una hora de incomodidad en la conversación
fundacional previene meses de conflicto irrecuperable más adelante.
\end{admonition}

La conexión entre esta sección y el material anterior del capítulo es estructural, no
incidental. Un fundador que modela la estimulación intelectual y la consideración individualizada
(liderazgo transformacional) mientras usa tácticas carismáticas para comunicar la visión (CLTs),
selecciona cofundadores por su aptitud cognitiva y responsabilidad, y estructura el entorno del
equipo para la seguridad psicológica ha ensamblado el kit de herramientas de liderazgo completo
que respalda la evidencia. Cada elemento refuerza a los demás: el carisma sin seguridad
psicológica produce obediencia del seguidor en lugar de aprendizaje; la seguridad psicológica
sin visión produce un equipo que genera ideas pero no puede priorizar; la visión sin la
disciplina interpersonal que \textcite{wasserman2012founders} describe produce un equipo fundador
que se fractura antes de que la visión haya tenido la oportunidad de ser probada. El capítulo de
comportamiento organizacional (\cref{ch:organizational-behavior}) traduce estos insumos de
liderazgo personal en los mecanismos estructurales --- diseño de incentivos, gobernanza y
claridad de roles --- que los sostienen a escala. \textcite{horowitz2014hard} es el relato profesional más citado sobre las decisiones de un CEO bajo presión --- los modos «en guerra» y «en paz», recortes de personal, despidos, y lo que el autor llama «la lucha» --- un complemento narrativo al conjunto de datos de 10{,}000 fundadores en \textcite{wasserman2012founders}.
